%=============================================================
%  REVISIÓN SISTEMÁTICA — Deep Learning para la detección de
%  anomalías en tráfico vehicular urbano y de autopista
%  Guías: Kitchenham y Charters (2007); PRISMA 2020; PRISMA-S
%  MATERIAL SUPLEMENTARIO de RSL_v3.tex (tablas por estudio)
%
%  Las cifras de resultados se generan con
%      python3 ../suplementos/generar_conteos.py
%  y se cargan desde generado/*.tex. No escribir conteos a mano.
%  Los datos que solo puede aportar el autor aparecen como
%  \pend{...} (en rojo en el PDF).
%=============================================================

\documentclass[12pt, a4paper]{article}

\usepackage[utf8]{inputenc}
\usepackage[T1]{fontenc}
\usepackage[spanish, es-nodecimaldot, es-tabla]{babel}
\usepackage{lmodern}
\usepackage{microtype}
\usepackage{setspace}
\onehalfspacing
\usepackage[top=3cm, bottom=2.5cm, left=3cm, right=2.5cm]{geometry}
\usepackage{amsmath, amssymb}
\usepackage{graphicx, float, caption}
\captionsetup{font=small, labelfont=bf}
\usepackage{tikz}
\usetikzlibrary{arrows.meta, positioning, calc}
\usepackage{booktabs, tabularx, array, multirow, longtable, xltabular}
\renewcommand{\arraystretch}{1.25}
\newcolumntype{Y}{>{\raggedright\arraybackslash}X}
\newcolumntype{P}[1]{>{\raggedright\arraybackslash}p{#1}}
\usepackage{pdflscape}
\usepackage{enumitem}
\setlist[itemize]{leftmargin=1.5em, itemsep=2pt}
\setlist[enumerate]{leftmargin=1.5em, itemsep=2pt}
\usepackage[table]{xcolor}
\definecolor{azulPrincipal}{RGB}{0, 70, 127}
\definecolor{azulMedio}{RGB}{0, 102, 179}
\definecolor{azulClaro}{RGB}{204, 229, 255}
\definecolor{grisClaro}{RGB}{245, 245, 245}
\definecolor{verdeOsc}{RGB}{0, 100, 60}
\definecolor{rojoAlerta}{RGB}{180, 0, 0}
\usepackage{tcolorbox}
\tcbuselibrary{skins, breakable}

\newtcolorbox{preguntabox}[1][]{enhanced,colback=azulClaro!50,colframe=azulPrincipal,
  fonttitle=\bfseries\color{white},coltitle=white,colbacktitle=azulPrincipal,
  title={Preguntas de investigaci\'{o}n},#1}
\newtcolorbox{hallazgo}[1][]{enhanced,breakable,colback=verdeOsc!8,colframe=verdeOsc,
  fonttitle=\bfseries\color{white},coltitle=white,colbacktitle=verdeOsc,
  title={Hallazgo},#1}

% \input primitivo: permite cargar filas dentro de tabular/longtable.
\makeatletter\newcommand{\tabinput}[1]{\@@input #1 }\makeatother

% Datos que debe completar el autor (aparecen en rojo).
\newcommand{\pend}[1]{\textcolor{rojoAlerta}{\textbf{[#1]}}}

\usepackage{fancyhdr}
\setlength{\headheight}{15pt}
\pagestyle{fancy}\fancyhf{}
\fancyhead[L]{\small\textcolor{azulPrincipal}{\textit{Material suplementario}}}
\fancyhead[R]{\small\textcolor{azulPrincipal}{\textit{Revisi\'{o}n sistem\'{a}tica}}}
\fancyfoot[C]{\small\thepage}
\renewcommand{\headrulewidth}{0.4pt}

\usepackage{titlesec}
\titleformat{\section}{\large\bfseries\color{azulPrincipal}}{\thesection.}{0.6em}{}[\titlerule]
\titleformat{\subsection}{\normalsize\bfseries\color{azulMedio}}{\thesubsection.}{0.5em}{}
\titleformat{\subsubsection}{\normalsize\itshape\color{azulPrincipal}}{\thesubsubsection.}{0.4em}{}

\usepackage{hyperref}
\hypersetup{colorlinks=true,linkcolor=azulPrincipal,citecolor=azulMedio,urlcolor=azulMedio,
  pdftitle={Material suplementario},
  pdfauthor={W. A. Mayta Arotaype}}
\usepackage{cite}

\renewcommand{\thetable}{S\arabic{table}}

%--- Cifras generadas desde la matriz de extracción ---------------------
% Generado por suplementos/generar_conteos.py. No editar a mano.
\newcommand{\Nestudios}{21}
\newcommand{\AnioMin}{2020}
\newcommand{\AnioMax}{2026}
\newcommand{\AnioA}{2}
\newcommand{\AnioB}{1}
\newcommand{\AnioC}{5}
\newcommand{\AnioD}{2}
\newcommand{\AnioE}{5}
\newcommand{\AnioF}{4}
\newcommand{\AnioG}{2}
\newcommand{\NrecientesVeintidosVeintiseis}{11}
\newcommand{\QAmedia}{4.38}
\newcommand{\QAde}{0.96}
\newcommand{\QAmin}{2.0}
\newcommand{\QAmax}{6.0}
\newcommand{\QAumbralAlta}{5.0}
\newcommand{\NaltaCalidad}{6}
\newcommand{\IDaltaCalidad}{S03, S04, S06, S10, S16, S17}
\newcommand{\QAUnosi}{6}
\newcommand{\QADossi}{7}
\newcommand{\QATressi}{12}
\newcommand{\QACuatrosi}{2}
\newcommand{\QACincosi}{2}
\newcommand{\QASeissi}{9}
\newcommand{\QASietesi}{11}
\newcommand{\QAOchosi}{6}
\newcommand{\Ncodigo}{13}
\newcommand{\NcodigoURL}{3}
\newcommand{\NetiqVerificadas}{6}
\newcommand{\NdatosSimuladosAny}{3}
\newcommand{\NdatasetPropio}{6}
\newcommand{\Nvariabilidad}{2}
\newcommand{\NdatoSensores}{4}
\newcommand{\NdatoMovilidad}{2}
\newcommand{\NdatoTrayectorias}{3}
\newcommand{\NdatoVideo}{7}
\newcommand{\NdatoMultimodal}{3}
\newcommand{\NdatoSimulado}{2}
\newcommand{\Nurbano}{7}
\newcommand{\Nautopista}{7}
\newcommand{\Nambos}{5}
\newcommand{\NentornoNE}{2}
\newcommand{\NfamRec}{14}
\newcommand{\NfamConv}{12}
\newcommand{\NfamAE}{6}
\newcommand{\NfamGrafo}{8}
\newcommand{\NfamTransf}{10}
\newcommand{\NfamGen}{2}
\newcommand{\Nhibrido}{17}
\newcommand{\Nmultimodal}{4}
\newcommand{\NetiqReal}{5}
\newcommand{\NetiqManual}{5}
\newcommand{\NetiqBench}{3}
\newcommand{\NetiqInyect}{4}
\newcommand{\NetiqMixto}{3}
\newcommand{\NetiqSin}{1}
\newcommand{\NetiqNR}{0}
\newcommand{\Ntasa}{9}
\newcommand{\Numbral}{19}
\newcommand{\Nparticion}{18}
\newcommand{\Nunidad}{20}
\newcommand{\Nmetaccuracy}{10}
\newcommand{\Nmetprecision}{7}
\newcommand{\Nmetrecall}{8}
\newcommand{\NmetFuno}{7}
\newcommand{\Nmetaucroc}{8}
\newcommand{\Nmetprauc}{2}
\newcommand{\Nmetdr}{4}
\newcommand{\Nmetfarfpr}{5}
\newcommand{\Nmetmttd}{2}
\newcommand{\Nmeterrormaermsemape}{4}
\newcommand{\Nmetioumap}{1}
\newcommand{\Nmetoperativaslatenciafpsenergia}{11}
\newcommand{\NcatA}{9}
\newcommand{\PcatA}{43~\%}
\newcommand{\NcatAnyA}{12}
\newcommand{\NcatB}{5}
\newcommand{\PcatB}{24~\%}
\newcommand{\NcatAnyB}{9}
\newcommand{\NcatC}{4}
\newcommand{\PcatC}{19~\%}
\newcommand{\NcatAnyC}{5}
\newcommand{\NcatD}{3}
\newcommand{\PcatD}{14~\%}
\newcommand{\NcatAnyD}{5}
\newcommand{\Nbaselines}{21}
\newcommand{\NbaselineNoDL}{11}
\newcommand{\NsinBaseline}{0}
\newcommand{\Nlimdatosreplicabilidad}{7}
\newcommand{\Nlimetiquetado}{9}
\newcommand{\Nlimgeneralizacion}{17}
\newcommand{\Nlimmetricasnouniformes}{1}
\newcommand{\Nlimcostodespliegue}{8}
\newcommand{\Nliminterpretabilidad}{1}
\newcommand{\XSensoresRec}{2}
\newcommand{\XSensoresConv}{1}
\newcommand{\XSensoresAE}{0}
\newcommand{\XSensoresGrafo}{2}
\newcommand{\XSensoresTransf}{2}
\newcommand{\XSensoresGen}{0}
\newcommand{\XMovilidadRec}{1}
\newcommand{\XMovilidadConv}{1}
\newcommand{\XMovilidadAE}{2}
\newcommand{\XMovilidadGrafo}{2}
\newcommand{\XMovilidadTransf}{0}
\newcommand{\XMovilidadGen}{0}
\newcommand{\XTrayectoriasRec}{3}
\newcommand{\XTrayectoriasConv}{2}
\newcommand{\XTrayectoriasAE}{1}
\newcommand{\XTrayectoriasGrafo}{2}
\newcommand{\XTrayectoriasTransf}{1}
\newcommand{\XTrayectoriasGen}{1}
\newcommand{\XVideoRec}{4}
\newcommand{\XVideoConv}{6}
\newcommand{\XVideoAE}{1}
\newcommand{\XVideoGrafo}{0}
\newcommand{\XVideoTransf}{2}
\newcommand{\XVideoGen}{0}
\newcommand{\XMultimodalRec}{3}
\newcommand{\XMultimodalConv}{2}
\newcommand{\XMultimodalAE}{1}
\newcommand{\XMultimodalGrafo}{1}
\newcommand{\XMultimodalTransf}{3}
\newcommand{\XMultimodalGen}{0}
\newcommand{\XSimuladoRec}{1}
\newcommand{\XSimuladoConv}{0}
\newcommand{\XSimuladoAE}{1}
\newcommand{\XSimuladoGrafo}{1}
\newcommand{\XSimuladoTransf}{2}
\newcommand{\XSimuladoGen}{1}
\newcommand{\NperiodoAntes}{3}
\newcommand{\XAntesRec}{1}
\newcommand{\XAntesConv}{3}
\newcommand{\XAntesAE}{0}
\newcommand{\XAntesGrafo}{1}
\newcommand{\XAntesTransf}{0}
\newcommand{\XAntesGen}{0}
\newcommand{\NperiodoMedio}{7}
\newcommand{\XMedioRec}{4}
\newcommand{\XMedioConv}{3}
\newcommand{\XMedioAE}{3}
\newcommand{\XMedioGrafo}{2}
\newcommand{\XMedioTransf}{2}
\newcommand{\XMedioGen}{1}
\newcommand{\NperiodoReciente}{11}
\newcommand{\XRecienteRec}{9}
\newcommand{\XRecienteConv}{6}
\newcommand{\XRecienteAE}{3}
\newcommand{\XRecienteGrafo}{5}
\newcommand{\XRecienteTransf}{8}
\newcommand{\XRecienteGen}{1}
\newcommand{\NgrafoNoVideo}{8}
\newcommand{\Nsupervisado}{10}
\newcommand{\Nnosupervisado}{11}
\newcommand{\NenfoquePred}{5}
\newcommand{\NenfoqueRecon}{5}
\newcommand{\NenfoqueClasif}{9}
\newcommand{\NetiqRealesTotal}{13}
\newcommand{\NetiqNoReales}{8}
\newcommand{\NumbralNoAplica}{8}
\newcommand{\NumbralExplicito}{11}
\newcommand{\NaltaEtiqReal}{3}
\newcommand{\NaltaHibrido}{5}
\newcommand{\NexcluidosLimite}{5}
\newcommand{\NincluidosLimite}{7}

% Generado por suplementos/prueba_sensibilidad_cadena.py. No editar a mano.
\newcommand{\NsensibOK}{19}
\newcommand{\IDsensibNo}{S03, S14}

%=============================================================
%  DATOS DE LA BÚSQUEDA Y DEL CRIBADO
%  Las fechas, los conteos por base, el flujo PRISMA y los kappa los
%  escribe automáticamente cribado/scripts/06_resultados_prisma.py.
%  A mano solo: \RegistroOSF y \RepositorioDatos.
%  Lo pendiente aparece en rojo en el PDF.
%=============================================================

% --- Fechas y registro ---------------------------------------
\newcommand{\FechaBusqueda}{19/10/2026}
\newcommand{\FechaBusquedaEN}{19 October 2026}
\newcommand{\RegistroOSF}{\pend{OSF, DOI 10.17605/OSF.IO/XXXXX}}
\newcommand{\RepositorioDatos}{\url{https://github.com/ClaudioB12/rsl-wilbert}}

% --- Revisor y fiabilidad del cribado (autor único) ------------
\newcommand{\Revisor}{W.\,A.\,M.\,A.}
\newcommand{\KappaRecribadoTA}{\pend{0,xx}}     % re-cribado ciego, título y resumen
\newcommand{\NRecribadoTA}{\pend{n}}
\newcommand{\KappaRecribadoTC}{\pend{0,xx}}     % re-cribado ciego, texto completo
\newcommand{\NRecribadoTC}{\pend{n}}
\newcommand{\KappaAsistenteTA}{0,06}     % acuerdo autor-asistente (TA)
\newcommand{\AcuerdoAsistenteTA}{31~\%}
\newcommand{\NExclRegla}{100}             % exclusiones propuestas por regla
\newcommand{\NExclReglaConfirmadas}{100}  % ... confirmadas por el autor
\newcommand{\KappaIATA}{0,77}            % acuerdo autor-IA en TA (pasa / no pasa)
\newcommand{\AcuerdoIATA}{93~\%}
\newcommand{\NCambiosIATA}{62}            % decisiones de la IA que el autor cambió
\newcommand{\NFueraFiltros}{123}           % quitados por filtros del protocolo aplicados por script
\newcommand{\NexclOA}{577}                 % excluidos en TA por no ser de acceso abierto (CI-2)
\newcommand{\FechaOA}{19/10/2026}             % fecha de verificación del acceso abierto
\newcommand{\AcuerdoCalidad}{\pend{xx~\%}}     % puntaje propuesto vs. verificado

% --- Registros por base (Tabla de bases) ----------------------
\newcommand{\NScopus}{498}
\newcommand{\NWoS}{193}
\newcommand{\NIEEE}{152}
\newcommand{\NScienceDirect}{337}
\newcommand{\NACM}{397}
\newcommand{\FScopus}{19/10/2026}
\newcommand{\FWoS}{19/10/2026}
\newcommand{\FIEEE}{19/10/2026}
\newcommand{\FScienceDirect}{19/10/2026}
\newcommand{\FACM}{19/10/2026}

% --- Flujo PRISMA 2020 ----------------------------------------
\newcommand{\Pidentificados}{1577}   % suma de las cinco bases
\newcommand{\Pduplicados}{337}
\newcommand{\Pcribados}{1240}
\newcommand{\PexclTA}{1137}
\newcommand{\Pbuscados}{103}        % solicitados a texto completo
\newcommand{\PnoRecuperados}{11}
\newcommand{\Pevaluados}{92}
\newcommand{\PexclTC}{20}
% Motivos de exclusión a texto completo (deben sumar \PexclTC)
\newcommand{\PexclCEdos}{2}      % CE-2 no es tráfico vehicular
\newcommand{\PexclCEtres}{7}     % CE-3 solo pronóstico/clasificación
\newcommand{\PexclCEcuatro}{2}   % CE-4 sin DL como método principal
\newcommand{\PexclCEcinco}{3}    % CE-5 tipo de documento / preprint
\newcommand{\PexclCIseis}{5}     % CI-6 sin evaluación cuantitativa de la detección
\newcommand{\PexclCEseis}{0}     % CE-6 información insuficiente
\newcommand{\PexclOtros}{1}      % CE-1, CI-2, CI-3
% Otros métodos (búsqueda v1 en ACM DL a texto completo)
\newcommand{\POtrosIdentificados}{0}
\newcommand{\POtrosIncluidos}{0}

% --- Base de recuperación de cada estudio (Anexo B) -----------
% La escribe cribado/scripts/06_resultados_prisma.py en generado/base_recuperacion.tex.
\newcommand{\BaseRecup}[1]{\ifcsname baserecup#1\endcsname\csname baserecup#1\endcsname\else\pend{base}\fi}

\IfFileExists{generado/base_recuperacion.tex}{% Generado por cribado/scripts/06_resultados_prisma.py. No editar a mano.
\expandafter\def\csname baserecupS01\endcsname{Scopus, Web of Science}
\expandafter\def\csname baserecupS02\endcsname{Scopus, Web of Science}
\expandafter\def\csname baserecupS03\endcsname{Scopus, Web of Science}
\expandafter\def\csname baserecupS04\endcsname{Scopus, Web of Science}
\expandafter\def\csname baserecupS05\endcsname{Scopus, Web of Science, IEEE Xplore}
\expandafter\def\csname baserecupS06\endcsname{ScienceDirect}
\expandafter\def\csname baserecupS07\endcsname{Web of Science}
\expandafter\def\csname baserecupS08\endcsname{ScienceDirect}
\expandafter\def\csname baserecupS09\endcsname{ScienceDirect}
\expandafter\def\csname baserecupS10\endcsname{Scopus, Web of Science, ScienceDirect}
\expandafter\def\csname baserecupS11\endcsname{Web of Science}
\expandafter\def\csname baserecupS12\endcsname{Scopus, Web of Science}
\expandafter\def\csname baserecupS13\endcsname{Scopus}
\expandafter\def\csname baserecupS14\endcsname{Scopus, ACM Digital Library}
\expandafter\def\csname baserecupS15\endcsname{Scopus, Web of Science}
\expandafter\def\csname baserecupS16\endcsname{ScienceDirect}
\expandafter\def\csname baserecupS17\endcsname{Scopus, Web of Science}
\expandafter\def\csname baserecupS18\endcsname{Scopus, ACM Digital Library}
\expandafter\def\csname baserecupS19\endcsname{Scopus, Web of Science}
\expandafter\def\csname baserecupS20\endcsname{Scopus, Web of Science}
\expandafter\def\csname baserecupS21\endcsname{Scopus, Web of Science}
\expandafter\def\csname baserecupS22\endcsname{Scopus}
\expandafter\def\csname baserecupS23\endcsname{Scopus}
\expandafter\def\csname baserecupS24\endcsname{Scopus}
\expandafter\def\csname baserecupS25\endcsname{Scopus, Web of Science}
\expandafter\def\csname baserecupS26\endcsname{Scopus, Web of Science}
\expandafter\def\csname baserecupS27\endcsname{Scopus, Web of Science, IEEE Xplore}
\expandafter\def\csname baserecupS28\endcsname{Scopus}
\expandafter\def\csname baserecupS29\endcsname{ScienceDirect}
\expandafter\def\csname baserecupS30\endcsname{Scopus, Web of Science}
\expandafter\def\csname baserecupS31\endcsname{Scopus}
\expandafter\def\csname baserecupS32\endcsname{Scopus, Web of Science}
\expandafter\def\csname baserecupS33\endcsname{ACM Digital Library}
\expandafter\def\csname baserecupS34\endcsname{ACM Digital Library}
\expandafter\def\csname baserecupS35\endcsname{ScienceDirect}
\expandafter\def\csname baserecupS36\endcsname{Scopus, Web of Science}
\expandafter\def\csname baserecupS37\endcsname{ScienceDirect}
\expandafter\def\csname baserecupS38\endcsname{Scopus, ScienceDirect}
\expandafter\def\csname baserecupS39\endcsname{ScienceDirect}
\expandafter\def\csname baserecupS40\endcsname{Scopus, Web of Science}
\expandafter\def\csname baserecupS41\endcsname{Scopus}
\expandafter\def\csname baserecupS42\endcsname{ScienceDirect}
\expandafter\def\csname baserecupS43\endcsname{Scopus}
\expandafter\def\csname baserecupS44\endcsname{Scopus, Web of Science}
\expandafter\def\csname baserecupS45\endcsname{Scopus, Web of Science}
\expandafter\def\csname baserecupS46\endcsname{ScienceDirect}
\expandafter\def\csname baserecupS47\endcsname{Scopus, Web of Science, ScienceDirect}
\expandafter\def\csname baserecupS48\endcsname{Scopus}
\expandafter\def\csname baserecupS49\endcsname{ACM Digital Library}
\expandafter\def\csname baserecupS50\endcsname{ScienceDirect}
\expandafter\def\csname baserecupS51\endcsname{Scopus, Web of Science, IEEE Xplore}
\expandafter\def\csname baserecupS52\endcsname{Web of Science}
\expandafter\def\csname baserecupS53\endcsname{ScienceDirect}
\expandafter\def\csname baserecupS54\endcsname{ACM Digital Library}
\expandafter\def\csname baserecupS55\endcsname{Scopus, Web of Science}
\expandafter\def\csname baserecupS56\endcsname{Scopus}
\expandafter\def\csname baserecupS57\endcsname{Scopus}
\expandafter\def\csname baserecupS58\endcsname{ScienceDirect}
\expandafter\def\csname baserecupS59\endcsname{Scopus, Web of Science}
\expandafter\def\csname baserecupS60\endcsname{Scopus}
\expandafter\def\csname baserecupS61\endcsname{Scopus, Web of Science}
\expandafter\def\csname baserecupS62\endcsname{Scopus}
\expandafter\def\csname baserecupS63\endcsname{Scopus, Web of Science, IEEE Xplore}
\expandafter\def\csname baserecupS64\endcsname{Scopus, Web of Science}
\expandafter\def\csname baserecupS65\endcsname{Scopus, Web of Science}
\expandafter\def\csname baserecupS66\endcsname{Scopus, Web of Science}
\expandafter\def\csname baserecupS67\endcsname{Scopus, Web of Science, ScienceDirect}
\expandafter\def\csname baserecupS68\endcsname{Scopus, Web of Science}
\expandafter\def\csname baserecupS69\endcsname{Scopus, Web of Science}
\expandafter\def\csname baserecupS70\endcsname{ScienceDirect}
\expandafter\def\csname baserecupS71\endcsname{Scopus, Web of Science}
\expandafter\def\csname baserecupS72\endcsname{Scopus, Web of Science, ScienceDirect}
}{}

%=============================================================
\begin{document}

\begin{center}
  {\Large\bfseries\color{azulPrincipal} Material suplementario}\\[0.6em]
  {\normalsize\textit{Deep Learning} para la detección de anomalías en tráfico vehicular urbano y de autopista: revisión sistemática de arquitecturas, datos, etiquetas y protocolos de evaluación (2020--2026)}\\[0.6em]
  {\small Wilbert Alex Mayta Arotaype --- Universidad Peruana Unión}
\end{center}
\noindent\rule{\linewidth}{1pt}

Contenido: Texto~S1 (cadenas de búsqueda); Tabla~S1 (historial del protocolo); Tabla~S2 (características de los estudios); Tabla~S3 (calidad por ítem y estudio); Tabla~S4 (líneas base); Tabla~S5 (estudios incluidos); Tabla~S6 (informes excluidos a texto completo); Tabla~S7 (lista de verificación PRISMA 2020). Las referencias de este material tienen numeración propia.

\subsection*{Cadena de b\'{u}squeda por base de datos}

Las cadenas se reproducen literalmente. Los filtros de a\~{n}o, tipo de documento e idioma se aplicaron con las opciones de la interfaz en Web of Science y ScienceDirect, y con un script sobre los metadatos exportados en Scopus, IEEE Xplore y ACM (Tabla~4 del artículo principal).

{\footnotesize
\noindent\textbf{Scopus} (b\'{u}squeda avanzada)
\begin{verbatim}
TITLE-ABS-KEY ( ( anomal* OR abnormal* OR "outlier detection" OR "incident
 detection" OR "accident detection" OR "accident prediction" OR "accident
 anticipation" OR "road closure*" OR "dangerous driving" OR "near-accident*"
 OR "near-miss*" OR "extreme event*" ) AND ( "deep learning" OR "deep neural"
 OR "neural network*" OR lstm OR gru OR cnn OR convolutional OR autoencoder*
 OR transformer* OR "graph neural" OR "self-attention" OR "attention
 mechanism" ) AND ( "road traffic" OR "urban traffic" OR "vehicular traffic"
 OR "highway traffic" OR "freeway traffic" OR "expressway traffic" OR "traffic
 flow*" OR "traffic speed" OR "traffic volume*" OR "traffic data" OR "traffic
 monitoring" OR "traffic video*" OR "traffic surveillance" OR "traffic
 camera*" OR "traffic scene*" OR "traffic scenario*" OR "traffic event*" OR
 "traffic incident*" OR "traffic accident*" OR "traffic anomal*" OR "traffic
 congestion" OR "road accident*" OR "road closure*" OR "intelligent
 transportation" OR "transportation network*" OR "road network*" OR "vehicle
 trajector*" OR "vehicle behavio*" OR "driving behavio*" OR taxi* ) ) AND
 PUBYEAR > 2019 AND PUBYEAR < 2027 AND ( LIMIT-TO ( DOCTYPE , "ar" ) OR
 LIMIT-TO ( DOCTYPE , "cp" ) ) AND ( LIMIT-TO ( LANGUAGE , "English" ) OR
 LIMIT-TO ( LANGUAGE , "Spanish" ) )
\end{verbatim}

\noindent\textbf{Web of Science Core Collection} (TS = t\'{i}tulo, resumen, palabras clave y Keywords Plus)
\begin{verbatim}
TS=((anomal* OR abnormal* OR "outlier detection" OR "incident detection" OR
 "accident detection" OR "accident prediction" OR "accident anticipation" OR
 "road closure*" OR "dangerous driving" OR "near-accident*" OR "near-miss*" OR
 "extreme event*") AND ("deep learning" OR "deep neural" OR "neural network*"
 OR LSTM OR GRU OR CNN OR convolutional OR autoencoder* OR transformer* OR
 "graph neural" OR "self-attention" OR "attention mechanism") AND ("road
 traffic" OR "urban traffic" OR "vehicular traffic" OR "highway traffic" OR
 "freeway traffic" OR "expressway traffic" OR "traffic flow*" OR "traffic
 speed" OR "traffic volume*" OR "traffic data" OR "traffic monitoring" OR
 "traffic video*" OR "traffic surveillance" OR "traffic camera*" OR "traffic
 scene*" OR "traffic scenario*" OR "traffic event*" OR "traffic incident*" OR
 "traffic accident*" OR "traffic anomal*" OR "traffic congestion" OR "road
 accident*" OR "road closure*" OR "intelligent transportation" OR
 "transportation network*" OR "road network*" OR "vehicle trajector*" OR
 "vehicle behavio*" OR "driving behavio*" OR taxi*))
Filtros (panel izquierdo): Publication Years 2020-2026; Document Types:
 Article, Proceeding Paper; Languages: English, Spanish.
\end{verbatim}

\noindent\textbf{IEEE Xplore} (Command Search, \textit{All Metadata}; m\'{a}ximo 25 t\'{e}rminos por consulta)
\begin{verbatim}
Consulta 1: ("All Metadata":anomal* OR "All Metadata":abnormal* OR "All
 Metadata":"incident detection" OR "All Metadata":"accident detection" OR "All
 Metadata":"accident prediction" OR "All Metadata":"road closure" OR "All
 Metadata":"dangerous driving" OR "All Metadata":"near-miss" OR "All
 Metadata":"extreme event") AND ("All Metadata":"deep learning" OR "All
 Metadata":"neural network" OR "All Metadata":LSTM OR "All Metadata":CNN OR
 "All Metadata":convolutional OR "All Metadata":autoencoder OR "All
 Metadata":transformer OR "All Metadata":"graph neural") AND ("All
 Metadata":"road traffic" OR "All Metadata":"urban traffic" OR "All
 Metadata":"highway traffic" OR "All Metadata":"traffic flow" OR "All
 Metadata":"traffic video" OR "All Metadata":"traffic surveillance" OR "All
 Metadata":"traffic incident" OR "All Metadata":"traffic accident")
Consulta 2: ("All Metadata":anomal* OR "All Metadata":abnormal* OR "All
 Metadata":"incident detection" OR "All Metadata":"accident detection" OR "All
 Metadata":"accident prediction" OR "All Metadata":"road closure" OR "All
 Metadata":"dangerous driving" OR "All Metadata":"near-miss" OR "All
 Metadata":"extreme event") AND ("All Metadata":"deep learning" OR "All
 Metadata":"neural network" OR "All Metadata":LSTM OR "All Metadata":CNN OR
 "All Metadata":convolutional OR "All Metadata":autoencoder OR "All
 Metadata":transformer OR "All Metadata":"graph neural") AND ("All
 Metadata":"traffic anomaly" OR "All Metadata":"traffic congestion" OR "All
 Metadata":"traffic data" OR "All Metadata":"traffic monitoring" OR "All
 Metadata":"intelligent transportation" OR "All Metadata":"road network" OR
 "All Metadata":"vehicle trajectory" OR "All Metadata":"driving behavior")
Consulta 3: ("All Metadata":anomal* OR "All Metadata":abnormal* OR "All
 Metadata":"incident detection" OR "All Metadata":"accident detection" OR "All
 Metadata":"accident prediction" OR "All Metadata":"road closure" OR "All
 Metadata":"dangerous driving" OR "All Metadata":"near-miss" OR "All
 Metadata":"extreme event") AND ("All Metadata":"deep learning" OR "All
 Metadata":"neural network" OR "All Metadata":LSTM OR "All Metadata":CNN OR
 "All Metadata":convolutional OR "All Metadata":autoencoder OR "All
 Metadata":transformer OR "All Metadata":"graph neural") AND ("All
 Metadata":"vehicular traffic" OR "All Metadata":"freeway traffic" OR "All
 Metadata":"traffic speed" OR "All Metadata":"traffic volume" OR "All
 Metadata":"traffic scene" OR "All Metadata":"traffic event" OR "All
 Metadata":"road accident" OR "All Metadata":taxi)
Filtros: Year 2020-2026; Content Type: Journals, Conferences.
\end{verbatim}

\noindent\textbf{ScienceDirect} (t\'{i}tulo, resumen y palabras clave; m\'{a}ximo 8 operadores, sin comodines)
\begin{verbatim}
Consulta 1: ("anomaly detection" OR "incident detection" OR "accident
 detection") AND ("deep learning" OR "neural network") AND ("road traffic" OR
 "urban traffic" OR "traffic flow")
Consulta 2: ("anomaly detection" OR "incident detection" OR "accident
 detection") AND ("deep learning" OR "neural network") AND ("traffic video" OR
 "traffic surveillance" OR "traffic accident")
Consulta 3: ("anomaly detection" OR "incident detection" OR "accident
 detection") AND ("deep learning" OR "neural network") AND ("traffic incident"
 OR "traffic congestion" OR "intelligent transportation")
Consulta 4: ("anomaly detection" OR "incident detection" OR "accident
 detection") AND ("deep learning" OR "neural network") AND ("vehicle
 trajectory" OR "road network" OR "highway traffic")
Consulta 5: ("abnormal" OR "anomalous") AND ("deep learning" OR "neural
 network") AND ("road traffic" OR "traffic flow" OR "vehicle trajectory")
Filtros: Years 2020-2026; Article type: Research articles.
\end{verbatim}

\noindent\textbf{ACM Digital Library} (campo Abstract)
\begin{verbatim}
[[Abstract: anomal*] OR [Abstract: abnormal*] OR [Abstract: "outlier
 detection"] OR [Abstract: "incident detection"] OR [Abstract: "accident
 detection"] OR [Abstract: "accident prediction"] OR [Abstract: "road
 closure*"] OR [Abstract: "dangerous driving"] OR [Abstract: "near-accident*"]
 OR [Abstract: "near-miss*"] OR [Abstract: "extreme event*"]] AND [[Abstract:
 "deep learning"] OR [Abstract: "deep neural"] OR [Abstract: "neural
 network*"] OR [Abstract: lstm] OR [Abstract: gru] OR [Abstract: cnn] OR
 [Abstract: convolutional] OR [Abstract: autoencoder*] OR [Abstract:
 transformer*] OR [Abstract: "graph neural"] OR [Abstract: "attention
 mechanism"]] AND [[Abstract: "road traffic"] OR [Abstract: "urban traffic"]
 OR [Abstract: "vehicular traffic"] OR [Abstract: "highway traffic"] OR
 [Abstract: "freeway traffic"] OR [Abstract: "traffic flow*"] OR [Abstract:
 "traffic speed"] OR [Abstract: "traffic volume*"] OR [Abstract: "traffic
 data"] OR [Abstract: "traffic monitoring"] OR [Abstract: "traffic video*"] OR
 [Abstract: "traffic surveillance"] OR [Abstract: "traffic camera*"] OR
 [Abstract: "traffic scene*"] OR [Abstract: "traffic event*"] OR [Abstract:
 "traffic incident*"] OR [Abstract: "traffic accident*"] OR [Abstract:
 "traffic anomal*"] OR [Abstract: "traffic congestion"] OR [Abstract: "road
 accident*"] OR [Abstract: "intelligent transportation"] OR [Abstract: "road
 network*"] OR [Abstract: "vehicle trajector*"] OR [Abstract: "driving
 behavio*"] OR [Abstract: taxi*]] AND [E-Publication Date: (01/01/2020 TO
 12/31/2026)]
Filtros: Content type: Research Article.
\end{verbatim}
}


\section*{Historial del protocolo}

\begin{table}[H]
\centering
\small
\caption{Historial de versiones del protocolo.}
\label{tab:s_historial}
\begin{tabularx}{\linewidth}{llY}
\toprule
\textbf{Versión} & \textbf{Fecha} & \textbf{Cambio} \\
\midrule
V1.0 & \pend{fecha} & Definición del tema, preguntas de investigación y PICOC. \\
V1.1 & \pend{fecha} & Cadenas por base (seis fuentes, incluidas MDPI y Taylor \& Francis). \\
V1.2 & \pend{fecha} & Cribado y eliminación de duplicados en Parsif.al. \\
V1.3 & \pend{fecha} & Evaluación de calidad (instrumento de 6 ítems) y extracción; 25 estudios. \\
V2.0 & \pend{fecha} & Respuesta a la revisión: alcance ampliado a video y autopistas; cinco bases (se añade IEEE Xplore y se retiran MDPI y Taylor \& Francis por ser editoriales); cadena validada con un conjunto de control; regla general para \textit{preprints} (CE-5); CI-5 y CI-6 precisados; nuevo instrumento de calidad de 8 ítems; nueva RQ6. \\
V2.1 & \pend{fecha} & Bloque C de la cadena acotado a frases ancladas al tráfico vial para reducir el ruido; la recuperación del conjunto de control se mantiene (19 de 21). \\
V3.0 & \FechaBusqueda & Búsqueda definitiva en las cinco bases; acceso abierto verificado registro por registro con OpenAlex y las bases (CI-2); cribado por título y resumen con propuestas de IA verificadas por el revisor; extracción y calidad a texto completo con propuestas citadas y verificadas; auditoría de coherencia (cinco inclusiones reclasificadas); corpus derivado solo de la búsqueda, sin otros métodos; la calidad deja de usarse para excluir y se usa para estratificar (alta $\geq$\,6, baja $<$\,4) con dos análisis de sensibilidad; re-cribado ciego intra-revisor; \Nestudios{} estudios incluidos. \\
\bottomrule
\end{tabularx}
\end{table}

\section*{Características, calidad y líneas base de los estudios incluidos}

\begin{landscape}
\footnotesize
\setlength{\tabcolsep}{3pt}
\begin{longtable}{llcP{2.2cm}P{1.3cm}P{3.0cm}P{2.0cm}P{2.4cm}P{3.5cm}P{3.6cm}}
\caption{Características de los estudios incluidos.\label{tab:caracteristicas}}\\
\toprule
\textbf{ID} & \textbf{Ref.} & \textbf{Año} & \textbf{Tipo de dato} & \textbf{Entorno} & \textbf{Datasets} & \textbf{Etiquetas} & \textbf{Anomalía} & \textbf{Modelo} & \textbf{Resultado principal} \\
\midrule\endfirsthead
\toprule
\textbf{ID} & \textbf{Ref.} & \textbf{Año} & \textbf{Tipo de dato} & \textbf{Entorno} & \textbf{Datasets} & \textbf{Etiquetas} & \textbf{Anomalía} & \textbf{Modelo} & \textbf{Resultado principal} \\
\midrule\endhead
\bottomrule\endfoot
\tabinput{generado/caracteristicas.tex}
\end{longtable}
\end{landscape}

\begin{small}
\begin{longtable}{lccccccccc}
\caption{Evaluación de calidad por ítem y estudio. Niveles: alta $\geq$\,\QAumbralAlta, baja $<$\,\QAumbralBaja (1 = Sí, 0.5 = Parcial, 0 = No).\label{tab:calidad}}\\
\toprule
\textbf{ID} & \textbf{QA1} & \textbf{QA2} & \textbf{QA3} & \textbf{QA4} & \textbf{QA5} & \textbf{QA6} & \textbf{QA7} & \textbf{QA8} & \textbf{Total} \\
\midrule\endfirsthead
\toprule
\textbf{ID} & \textbf{QA1} & \textbf{QA2} & \textbf{QA3} & \textbf{QA4} & \textbf{QA5} & \textbf{QA6} & \textbf{QA7} & \textbf{QA8} & \textbf{Total} \\
\midrule\endhead
\bottomrule\endfoot
\tabinput{generado/calidad.tex}
\end{longtable}
\end{small}

\begin{landscape}
\footnotesize
\begin{longtable}{lP{7.5cm}cP{8.5cm}}
\caption{Líneas base y ventaja frente a la mejor de ellas (RQ6).\label{tab:baselines}}\\
\toprule
\textbf{ID} & \textbf{Líneas base} & \textbf{No DL} & \textbf{Ventaja frente a la mejor línea base} \\
\midrule\endfirsthead
\toprule
\textbf{ID} & \textbf{Líneas base} & \textbf{No DL} & \textbf{Ventaja frente a la mejor línea base} \\
\midrule\endhead
\bottomrule\endfoot
\tabinput{generado/baselines.tex}
\end{longtable}
\end{landscape}

\section*{Estudios incluidos y excluidos}

\begin{landscape}
\footnotesize
\begin{longtable}{lllP{2.8cm}cP{3.0cm}P{5.5cm}P{3.0cm}}
\caption{Estudios incluidos con su identificador de cribado, DOI y base de recuperación.\label{tab:ap_estudios}}\\
\toprule
\textbf{ID} & \textbf{ID cribado} & \textbf{Ref.} & \textbf{Autores} & \textbf{Año} & \textbf{Revista o conferencia} & \textbf{DOI} & \textbf{Base de recuperación} \\
\midrule\endfirsthead
\toprule
\textbf{ID} & \textbf{ID cribado} & \textbf{Ref.} & \textbf{Autores} & \textbf{Año} & \textbf{Revista o conferencia} & \textbf{DOI} & \textbf{Base de recuperación} \\
\midrule\endhead
\bottomrule\endfoot
\tabinput{generado/lista_estudios.tex}
\end{longtable}
\end{landscape}

\begin{small}
\begin{longtable}{P{4.2cm}P{2.0cm}P{8.2cm}}
\caption{Informes evaluados a texto completo y excluidos, con el criterio aplicado y su justificación (PRISMA, ítem 16b).\label{tab:limite}}\\
\toprule
\textbf{Informe} & \textbf{Decisión} & \textbf{Justificación} \\
\midrule\endfirsthead
\toprule
\textbf{Informe} & \textbf{Decisión} & \textbf{Justificación} \\
\midrule\endhead
\bottomrule\endfoot
\tabinput{generado/limite.tex}
\end{longtable}
\end{small}

\section*{Lista de verificación PRISMA 2020}

\begin{small}
\begin{longtable}{lP{8.8cm}P{4.0cm}}
\caption{Ubicación de los ítems de PRISMA 2020. Las secciones, tablas y figuras sin prefijo «S» corresponden al artículo principal.\label{tab:prisma_check}}\\
\toprule
\textbf{Ítem} & \textbf{Requisito} & \textbf{Ubicación} \\
\midrule\endfirsthead
\toprule
\textbf{Ítem} & \textbf{Requisito} & \textbf{Ubicación} \\
\midrule\endhead
\bottomrule\endfoot
1 & Título que identifica el trabajo como revisión sistemática & Título \\
2 & Resumen estructurado (PRISMA for Abstracts) & Resumen \\
3 & Justificación en relación con el conocimiento existente & Sec.~1, 1.1; Tabla~1 \\
4 & Objetivos o preguntas & Sec.~1 (RQ1--RQ6) \\
5 & Criterios de elegibilidad & Tabla~5 \\
6 & Fuentes de información y fecha de última búsqueda & Sec.~3.3; Tabla~4 \\
7 & Estrategia de búsqueda completa & Texto~S1 \\
8 & Proceso de selección & Sec.~3.5 \\
9 & Proceso de extracción & Sec.~3.5 y Sec.~3 (Extracción de datos) \\
10 & Variables extraídas & Tabla~6 \\
11 & Evaluación del riesgo de sesgo & Sec.~3; Tabla~7 \\
12 & Medidas de efecto & No aplica (síntesis sin metaanálisis); Sec.~3.8 \\
13a--e & Métodos de síntesis & Sec.~3.8 \\
13f & Análisis de sensibilidad & Sec.~3.8 y 4.10 \\
14 & Sesgo de reporte & Sec.~3.8, justificación \\
15 & Certeza de la evidencia & Sec.~3.8, justificación \\
16a & Resultados de la búsqueda y selección & Figura~1 \\
16b & Estudios excluidos que parecían cumplir los criterios & Tabla~S6 \\
17 & Características de los estudios & Tabla~S2; Tabla~S5 \\
18 & Riesgo de sesgo por estudio & Tabla~S3 \\
19 & Resultados de cada estudio & Tablas~S2 y S4 \\
20a--d & Resultados de las síntesis & Sec.~4 (RQ1--RQ6); Tabla~18 \\
21 & Sesgo de reporte en la síntesis & Sec.~4.9 \\
22 & Certeza de la evidencia & Sec.~5.4 \\
23a--d & Discusión, limitaciones e implicaciones & Sec.~5 \\
24a--c & Registro y protocolo & Declaraciones; Texto~S1 \\
25 & Financiamiento & Declaraciones \\
26 & Conflicto de intereses & Declaraciones \\
27 & Disponibilidad de datos, código y materiales & Declaraciones \\
\end{longtable}
\end{small}

\bibliographystyle{IEEEtran}
\bibliography{referencias,estudios}

\end{document}
